% ---------------------------------------------------------------------
% 2.9 Redes convolucionales y recurrentes para la detección de la P300
% ---------------------------------------------------------------------
\section[REDES CONVOLUCIONALES Y RECURRENTES PARA LA DETECCIÓN DE LA P300]{Redes convolucionales
y recurrentes para la detección de la P300}
\label{sec:cnn}

\subsection{Redes neuronales convolucionales}

Las redes neuronales convolucionales (CNN) están diseñadas para trabajar con datos organizados en
forma de cuadrícula, como las imágenes o las series de tiempo, en los que los valores cercanos
están muy relacionados. Su nombre se debe a que usan la operación de convolución en al menos una de
sus capas \cite{goodfellow2016}. En una capa
convolucional, un filtro o \textit{kernel} de pesos aprendibles se desplaza sobre la entrada y en
cada posición calcula la suma ponderada de los valores que cubre, de forma análoga a la
Ecuación~\eqref{eq:convolucion}. El resultado se denomina mapa de características. Como el mismo
filtro se aplica en todas las posiciones, el número de parámetros es mucho menor que en una capa
totalmente conectada \cite{aggarwal2018}. Además, estas capas tienen tres parámetros
importantes \cite{goodfellow2016}:
\begin{itemize}
  \item \textbf{Paso (\textit{stride}):} controla cuánto se desplaza el filtro en cada posición.
  \item \textbf{Relleno (\textit{padding}):} agrega ceros en los bordes para conservar la longitud de la salida.
  \item \textbf{Agrupamiento (\textit{pooling}):} reduce la resolución de los mapas conservando el valor máximo o promedio de cada región.
\end{itemize}

En el caso del EEG, cada época puede verse como una matriz de $C$ canales por $T$ muestras, lo que
permite separar el procesamiento en dos dimensiones con significado fisiológico distinto: las
convoluciones temporales, con filtros de tamaño $1\times K$, aprenden patrones en el tiempo
equivalentes a filtros de frecuencia, mientras que las convoluciones espaciales, con filtros de
tamaño $C\times1$, aprenden combinaciones ponderadas de los electrodos equivalentes a un filtrado
espacial. Cecotti y Gräser aplicaron esta separación en la primera CNN diseñada específicamente para
detectar la P300, en la que una primera capa espacial es seguida de una capa temporal, alcanzando
en la base de datos de la competencia BCI III una exactitud de alrededor del 95\% con 15 repeticiones
\cite{cecotti2011}. Esta misma arquitectura fue implementada en el trabajo preliminar.

\subsection{Convoluciones separables en profundidad}

Una convolución separable en profundidad descompone la convolución estándar en dos pasos: primero
una convolución \textit{depthwise}, que aplica un filtro independiente a cada canal de entrada, y
después una convolución \textit{pointwise} de tamaño $1\times1$, que combina los canales
resultantes \cite{chollet2017}. Esta descomposición reduce drásticamente el número de parámetros respecto a una convolución
estándar \cite{chollet2017}. Esta idea es la base de dos de las arquitecturas más compactas para EEG,
EEGNet \cite{lawhern2018} y SepConv1D \cite{alvarado2021}, que resultan especialmente convenientes
cuando se dispone de pocos datos o cuando el modelo debe ejecutarse en dispositivos con recursos
limitados.

\subsection{EEGNet}

EEGNet es una red convolucional compacta propuesta por Lawhern et al. para clasificar señales de
EEG en distintos paradigmas de BCI, entre ellos la P300, con un número muy reducido de parámetros
\cite{lawhern2018}. Como se muestra en la Figura~\ref{fig:eegnet}, su primer bloque aplica una
convolución temporal que funciona como un banco de filtros de frecuencia. Después, una convolución
\textit{depthwise} espacial aprende, para cada filtro temporal, una combinación de electrodos, lo
que equivale a un filtrado espacial. El segundo bloque aplica una convolución separable que resume
cada mapa de características en el tiempo y combina los mapas entre sí. Por último, la
clasificación se realiza con una capa densa con activación softmax \cite{lawhern2018}. En el trabajo preliminar esta arquitectura se implementó con ocho filtros
temporales de tamaño $1\times16$ sobre épocas de 64 canales por 78 muestras.

\begin{figure}[htbp]
\centering
\begin{tikzpicture}[
  capa/.style={draw, rounded corners=2pt, align=center, minimum height=1.35cm, text width=2.0cm, font=\scriptsize},
  flecha/.style={-{Stealth[length=1.8mm]}, thick}, node distance=0.28cm]
  \node[capa, fill=gray!8] (in) {Entrada\\$C\times T$\\(64 $\times$ 78)};
  \node[capa, fill=blue!8, right=of in] (c1) {Conv. temporal\\$F_1$ filtros $1\times K$\\+ BatchNorm};
  \node[capa, fill=blue!8, right=of c1] (dw) {Conv. \textit{depthwise}\\espacial $C\times1$\\BN + ELU\\Pool + Dropout};
  \node[capa, fill=red!8, right=of dw] (sep) {Conv. separable\\$F_2$ filtros\\BN + ELU\\Pool + Dropout};
  \node[capa, fill=gray!8, right=of sep, text width=1.6cm] (out) {Densa\\+ softmax\\P300 / no P300};
  \foreach \a/\b in {in/c1, c1/dw, dw/sep, sep/out} \draw[flecha] (\a) -- (\b);
  \draw[decorate, decoration={brace, amplitude=5pt, mirror}] ([yshift=-3pt]c1.south west) -- node[below=6pt, font=\scriptsize] {Bloque 1: filtrado temporal y espacial} ([yshift=-3pt]dw.south east);
  \draw[decorate, decoration={brace, amplitude=5pt, mirror}] ([yshift=-3pt]sep.south west) -- node[below=6pt, font=\scriptsize] {Bloque 2} ([yshift=-3pt]sep.south east);
\end{tikzpicture}
\caption{Arquitectura general de EEGNet. Elaboración propia con base en \cite{lawhern2018}.}
\label{fig:eegnet}
\end{figure}

\subsection{Arquitecturas convolucionales para la detección de la P300}

A partir del trabajo de Cecotti y Gräser se han propuesto numerosas arquitecturas convolucionales
para detectar la P300. Alvarado-González et al. compararon doce de ellas en cuatro bases de datos
distintas, bajo los esquemas intrasujeto y entre sujetos, y encontraron que las diferencias de AUC
entre arquitecturas eran pequeñas, al grado de que SepConv1D, con una sola capa convolucional
separable de cuatro filtros seguida de una neurona sigmoide, obtuvo un desempeño comparable al de
redes con muchos más parámetros \cite{alvarado2021}. La Tabla~\ref{tab:arquitecturas} resume algunas de
las arquitecturas evaluadas y los valores de AUC reportados. Cabe mencionar que Contreras utilizó
SepConv1D en su sistema con cuatro canales \cite{contreras2021}.

\begin{table}[htbp]
\centering
\caption{Algunas arquitecturas convolucionales para la detección de la P300 y AUC promedio reportada en
\cite{alvarado2021}}
\label{tab:arquitecturas}
\small
\begin{tabularx}{\textwidth}{L{2.8cm}YL{1.3cm}L{1.3cm}}
\toprule
\textbf{Arquitectura} & \textbf{Descripción} & \textbf{AUC intra} & \textbf{AUC inter} \\
\midrule
CNN-1 \cite{cecotti2011} & 4 capas: convolución espacial (10 filtros) y temporal con tangente hiperbólica escalada; salida sigmoide. & 0.82 & 0.78 \\
DeepConvNet \cite{schirrmeister2017} & Cuatro bloques convolucionales seguidos de softmax. & 0.84 & 0.79 \\
EEGNet \cite{lawhern2018} & Red compacta con convoluciones \textit{depthwise} y separables. & 0.84 & 0.80 \\
SepConv1D \cite{alvarado2021} & Una capa convolucional 1D separable (4 filtros) y una neurona sigmoide. & 0.84 & 0.78 \\
\bottomrule
\end{tabularx}
\par\smallskip{\footnotesize Elaboración propia con información de \cite{alvarado2021,contreras2021}.}
\end{table}

\subsection{Redes neuronales recurrentes}

A diferencia de las redes convolucionales, las redes neuronales recurrentes (RNN) procesan una
secuencia elemento por elemento y mantienen un estado oculto que resume la información de los
pasos anteriores, lo que las hace adecuadas para modelar series de tiempo \cite{goodfellow2016}.
Su principal limitación es el desvanecimiento del gradiente, que les impide aprender dependencias
de largo plazo. Para resolverlo, Hochreiter y Schmidhuber propusieron las redes de memoria de
corto y largo plazo (LSTM). Estas incorporan una celda de memoria y tres compuertas (de entrada, de
olvido y de salida), que regulan qué información se agrega, se descarta o se entrega en cada paso
\cite{hochreiter1997}. En el contexto de la P300 se han utilizado principalmente en arquitecturas
híbridas, en las que una CNN extrae características de cada época y una capa recurrente modela su
evolución temporal, enfoque que también fue explorado en el trabajo preliminar.

